\section*{Hoofdstuk \ref{ch:g4:fractions} --- Eerste breuken}

\begin{solution}{exo:g4:fractions:1}
Pizza: $\frac56$ is over. Chocolade: $\frac38$ opgegeten. Vlag: $\frac13$
gekleurd.
\end{solution}

\begin{solution}{exo:g4:fractions:2}
$\frac35$ kleurt drie delen in, $\frac25$ maar twee: $\frac35$ geeft het
grootste gekleurde stuk.
\end{solution}

\begin{solution}{exo:g4:fractions:3}
\omterm{def:g3:division:half}{Helft} van $18$: $9$. \omterm{def:g3:division:half}{Kwart} van $20$: $5$. Derde van $21$: $7$. $\frac34$
van $20$: drie \omterm{def:g3:division:half}{kwarten}, $3 \times 5 = 15$.
\end{solution}

\begin{solution}{exo:g4:fractions:4}
$\frac13$: één stapje; $\frac33$: op $1$; $\frac53$: twee stapjes voorbij
$1$. De \omterm{def:g4:fractions:def}{breuk} die precies op $2$ uitkomt, is $\frac63$.
\end{solution}

\begin{solution}{exo:g4:fractions:5}
$\frac38 < \frac68$; \quad $\frac55 = 1$; \quad $\frac74 > 1$;
\quad $\frac12 = \frac24$.
\end{solution}

\begin{solution}{exo:g4:fractions:6}
$\frac26 + \frac36 = \frac56$; \quad
$\frac58 - \frac28 = \frac38$; \quad
$\frac14 + \frac24 = \frac34$; \quad
$1 = \frac33$, dus $1 - \frac13 = \frac23$.
\end{solution}

\begin{solution}{exo:g4:fractions:7}
$\frac54 = 1 + \frac14$; \quad
$\frac73 = 2 + \frac13$; \quad
$\frac92 = 4 + \frac12$.
\end{solution}

\begin{solution}{exo:g4:fractions:8}
Een half uur: $30$ min. Een \omterm{def:g3:division:half}{kwart}: $15$ min. Driekwart: $45$ min. Een
derde: $20$ min.
\end{solution}

\begin{solution}{exo:g4:fractions:9}
$\frac14$ van $24$ is $6$ leerlingen met een bril; $24 - 6 = 18$ zonder.
\end{solution}

\begin{solution}{exo:g4:fractions:10}
Samen: $\frac38 + \frac48 = \frac78$. Over: $\frac18$. Sam heeft er meer
gegeten ($4$ achtsten tegen $3$).
\end{solution}

\begin{solution}{exo:g4:fractions:11}
Dat kan inderdaad: de \omterm{def:g3:division:half}{helft} van een kleine pizza kan minder eten zijn dan
een \omterm{def:g3:division:half}{kwart} van een reuzenpizza. \omterm{def:g4:fractions:def}{Breuken} vergelijken delen \emph{van hetzelfde
geheel}; voordat je $\frac12$ en $\frac14$ ``van iets'' vergelijkt, moet je
weten dat die twee gehelen even groot zijn.
\end{solution}
